The field of research on code generation with artificial intelligence has been growing exponentially in recent years. Both academia and industry have produced numerous proposals, tools, and models to support this activity, making it one of the most dynamic topics in contemporary software engineering. This scenario highlights a rapidly evolving field that necessitates new research capable of providing timely, comprehensive reviews of available code-generation options and their specific details. Given the breadth of recent publications, for the sake of objectivity, this article highlights only a few works directly relevant to the topic, without attempting to exhaust the entire body of literature in the area.

\cite{jiang2024survey} presents a systematic and comprehensive review of the state of the art in language models applied to code generation. In addition to proposing a taxonomy that covers data curation, technical advances, and even ethical and environmental impacts, the study also provides a chronological overview of the evolution of large language models for code generation in recent years, highlighting the rapid succession of innovations in this field. This work relates to the present article by demonstrating the consolidation of code generation as its own research line, with accelerated expansion, serving as a foundation for identifying advances and gaps that still require further exploration, especially regarding empirical evaluation and the gap between academia and industrial practice.

\cite{liu2024large} broadens the discussion by analyzing the use of LLM-based agents in the context of software engineering. In addition to positioning agents as a new paradigm—where models are combined with mechanisms of planning, memory, perception, and action to deal with complex tasks—the study provides a comprehensive overview of their application in different software engineering tasks, including code generation, and lists several agents and tools that have been employed for this purpose. This perspective is relevant because it shows a trend in the evolution of LLM usage, from simple code-generation assistants to autonomous systems capable of acting across end-to-end development processes.

\cite{sergeyuk2025using} offer a more applied perspective and are closer to the idea of the present article, as they investigate how industry is using artificial intelligence for code generation. However, this analysis was conducted through interviews with 481 developers, which enabled the mapping of perceptions, practices, and limitations of these assistants in everyday software development. The authors highlight which activities programmers prefer to delegate to AI, as well as the reasons why these assistants are not used in certain stages, such as a lack of trust or contextual limitations. This work complements previous work by providing practical evidence of adoption and the challenges faced in real-world development environments, offering a user-centered, practice-oriented perspective.

Finally, unlike the highlighted works, this article does not aim solely to systematize technical advances or gather usage perceptions. Instead, it distinguishes itself by proposing a rapid and multivocal review. By adopting this approach, the goal is to capture not only what the formal academic literature has discussed but also perspectives from technical reports, specialized blogs, and other industry sources, aiming to provide a broader and more up-to-date view of what is effectively being used in terms of artificial intelligence for code generation, thus accelerating the understanding of the current state of this constantly evolving field.

% Multivocal literature reviews have increasingly gained attention as methods to integrate insights from both academic literature and grey literature, providing richer and broader perspectives on technological topics. In software engineering and related areas, several studies have successfully demonstrated the value of multivocal reviews for understanding complex phenomena involving AI technologies. For example, Felizardo and Carver \cite{felizardo2020automating} outline automated systematic literature reviews in software engineering, emphasizing the importance of considering diverse information sources to ensure comprehensive analyses.

% Cartaxo et al. \cite{cartaxo2020rapid} discuss how rapid reviews can efficiently summarize evidence to support timely decision-making in dynamic fields such as software engineering. Moreover, Garousi et al. \cite{garousi2020grey} explicitly address the significance of incorporating grey literature into software engineering research, arguing that industrial and practitioner-driven insights enrich the applicability and context-relevance of findings.

% In the broader AI context, similar methodological frameworks have been applied. Previous work \cite{swbokv4} has emphasized the evolving role of generative models and large language models, underscoring the relevance of multivocal approaches to stay updated with rapidly evolving AI capabilities. Thus, this rapid multivocal review not only draws from these methodological precedents but also expands upon them by specifically examining the state-of-the-art and practice of generative AI solutions applicable to software construction.